\documentclass{article}
\usepackage[T1]{fontenc}
\RequirePackage[utf8]{inputenc}
\RequirePackage[french]{babel}
\RequirePackage[
    top=2.5cm, bottom=2.5cm,
    left=2.8cm, right=2.5cm,
    headheight=14pt
]{geometry}
\RequirePackage{fancyhdr}
\usepackage{graphicx} % Required for inserting images

\title{resume de la méthodologie d'investigation}
\author{AYONMENE TIOBOU Varese 22p045 HN4-CIN}
\date{March 2026}

\begin{document}

\maketitle

\section*{\begin{center}
 PROCESSUS INVESTIGATIF UNIVERSEL (PIU)   
\end{center}}
Le processus investigatif universel est une méthodologie utilise dans le cadre de tout investigation à savoir celle impliquant un cadre personnel (enquête personnel), professionnel (investigation interne a une organisation) et judiciaire (initie par une autorité juridictionnelle). Ce processus englobe 4 objectifs principale qui sont les suivantes :
\begin{itemize}
    \item Établissement factuel : les faits sont vérifiables et reproductibles
	\item Intégrité probatoire : les preuves ne sont pas altérées
    \item Défendabilité : résisté a toute contradictions possibles
	\item Traçabilité : toute les actions sont documentées
\end{itemize}
Ce processus se réalise en 5 phase à savoir : l’initialisation, l’acquisition, investigations, la formalisation et le suivi. Ces phases sont découpées en un ensemble de 17 étapes opérationnelles évaluer à l’aide de trois quality gate (point de contrôle).
\subsection*{PHASE 1 : L’INITIALISATION}
Cette phase regroupe les trois premières étapes du processus et permet d’établir les conditions opérationnelles de l’exécution de l’investigation.
\begin{itemize} 
	\item
    La première étape intitule « analyse de la demande » permet la réception et compréhension de la demande avant la prise de décision de tout engagements nécessaires qui ce clôture par la production d’une fiche d’évaluation initiale.
	\item La deuxième étape intitule « décision d’engagement » consiste en la validation de l’opération d’investigation d’après les critères de gouvernance propre au contexte et ce solde par la réalisation d’un mandat d’investigations
	\item La troisième étape intitule « planification opérationnelle » qui consiste a structure la mission tout en réunissant les ressources nécessaire et réalisé une projection dans le futur afin d’anticipé les difficultés opérationnelles ce solde par la livraison d’un plan d’investigation.
\end{itemize}
\subsection*{QUALITY GATE 1 : VALIDATION AVANT L’OPERATION D’INVESTIGATION}
Permet de s’assurer que l’ensemble des conditions préalables au démarrage de l’investigation sont satisfaites.
\subsection*{PHASE 2 : ACQUISITION}
Consiste en la collecte des preuves et des éléments nécessaires tout en garantissant l’intégrité probatoire. Elle est constituée des trois étapes suivantes celle de la phase 1.
\begin{itemize}
	\item Etape 4 : collecte Systématique des éléments probatoires qui consiste au réassemblage exhaustif des éléments factuel pertinents tout en gardant leur intégrité et leur valeur probatoire et ce solde par la production d’un inventaire des éléments probatoires
	\item Etape 5 : sécurisation et préservation permet de consolider l’intégrité des éléments collectées et éviter tout modifications, la perte ou un accès non autorise et permet de produire un registre de préservation.
	\item Etape 6 : identification et cartographie des parties prenantes qui se résulte par la production d’une matrice des parties prenantes
\end{itemize}
\subsection*{QUALITY GATE 2 : INTEGRITE PROBATOIRE}
Permet de s’assurer que la base probatoire constitue est complète, intègre, traçable et juridiquement défendable avant l’exécution de la phase analytique.
\subsection*{PHASE 3 : INVESTIGATION ANALYTIQUE}
Cette phase constitue le cœur du processus au cours de laquelle des hypothèses sont émise et confrontées aux éléments probatoires et elle regroupe les étapes 7 à 10.
\begin{itemize}
	\item Etape 7 : interrogatoires et collecte testimoniale qui consiste à la colleté des témoignages et expertises orales permettant la compréhension approfondie des faits et conclus à la production d’un corpus testimoniales.
	\item Etape 8 : génération d’hypothèses multiples qui consiste à établir des scenarios explicatifs sans présupposé initial pour éviter les biais et ce conclut par la production d’un arbre d’hypothèses.
	\item Etape 9 : confrontation hypothèses/preuves ou l’évaluation de la compatibilité de chaque hypothèse avec les preuves et les témoignages collectes. Elle aboutit à la production d’une matrice d’analyse d’hypothèses/preuves après quoi vient une falsification active qui conduit à la production d’un rapport de falsification.
	\item Etape 10 : validation et vérification qui consiste à la soumission des conclusions préliminaires émergentes à un processus de validations multiméthodes qui conduit à la production d’un rapport de validation.
\end{itemize}
\subsection*{QUALITY GATE 3 : ROBUSTESSE ANALYTIQUE}
S’assurer de la robustesse, de la documentation, et de la défensabilité des conclusions investigatrices pour la justification de leur formalisation et la transmission au commanditaire.
\subsection*{PHASE 4 : FOEMALISATION ET TRANSMISSION}
\begin{itemize}
     

	\item Etape 11 : documentation exhaustive et consolidation l’objectif étant de rassembler, structurer et archiver la totalité des éléments du dossier dans un format pérennes et auditable. Ce document doit comporter les parties suivantes : synthèse, analyses, preuves, témoignages et annexes. 
	\item Etape 12 : rédaction du livrable rapport qui consiste à produire un document de synthèse présentant les conclusions de l’investigation comportant la structure suivante : synthèse exécutive, rapport technique, dossier probatoire et ce document doit être clair, objective, défendable et adaptable.
	\item Etape 13 : transmission sécurisée du livrable ici il est question de transmettre le livrable au commanditaire de la manière la plus sécurisée et tracable possible conforme aux exigences de confidentialité et adaptée au contexte juridique. Elle doit ce solde par l’obtention d’une preuve de transmission
\end{itemize}
\subsection*{PHASE 5 : SUIVI POST-TRANSMISSION}
\begin{itemize}
     

	\item Etape 14 : préparation à la défense contradictoire en anticipant sur les contestations possibles et la préparation d’une défense rigoureuse et documentée des conclusions investigatrices
	\item Etape 15 : maintenance probatoire continue en maintenant l’intégrité et l’accessibilité du dossier durant toute la procédure contentieuse ou administrative.
	\item Etape 16 : analyse des contre-expertises via l’évaluation des rapport contradictoires produits par la partie adverse et la préparation d’une réfutation méthodique et professionnelle.
	\item Etape 17 : clôture et retour d’expériences par l’archivage définitive du dossier, la capitalisation de l’apprentissage et l’amélioration des processus futur. 
\end{itemize}
\end{document}
